\section{Preparation - Before measurements are performed}
Once a field of study has been defined and experimental data is to be collected, several preparatory aspects should be considered. The experimenter should aim to familiarise themselves with the physical nature of the phenomenon or quantity to be observed and, where possible, develop an intuitive understanding of its expected behaviour. It is acknowledged, however, that complex systems do not always exhibit behaviour that is immediately intuitive.

This section describes the key elements that should be prepared before final measurements are initiated. In practice, this process is often iterative, with preliminary measurements performed to gain familiarity with the instruments, the experimental setup, and the general behaviour of the system under investigation.

\subsection{Define the Quantity}

The first step in any measurement campaign is to clearly define the quantity or phenomenon to be measured or observed. Once this is established, a method for observing or quantifying that quantity must be designed.

In metrology, this is often referred to as an \emph{indirect measurement}, since in many cases we cannot directly observe the quantity of interest. A simple example is measuring electrical current. There are no methods to directly measure current in a straightforward way; instead, we measure a related, observable quantity. Common approaches include:

\begin{itemize}
\item Using a shunt resistor to measure the voltage drop across it.
\item Employing a Hall-effect sensor to measure the magnetic field generated by the current.
\item Measuring the magnetic flux through a coil induced by the current.
\end{itemize}

In each case, we ultimately measure a voltage that is related to the current. This is considered an indirect measurement, where the desired quantity (current) is inferred from a relationship to a directly observable quantity (voltage, magnetic field, or flux).

Although the previous example is simple, the principle of indirect measurement applies broadly. For instance, one might wish to observe the RF spectrum or the reflection coefficients that arise when a hand is moved near an antenna. In this case, the indirect measurement is the reflection coefficient, while the desired phenomenon is the placement or motion of the hand. By measuring the reflection coefficient, we infer the effect of the hand on the antenna, illustrating how indirect measurements allow us to quantify phenomena that cannot be directly observed.

It is important to note that in exploratory work - such as investigating whether the position of a hand can be inferred from measured reflection coefficients - the principle of indirect measurement still applies. In this case, we must not only measure the reflection coefficients, but also know the exact placement of the hand. This allows us to map a relationship between the observable quantity (reflection coefficients) and the underlying phenomenon (hand position), which is essential for drawing meaningful conclusions.

It is also acknowledged that some experimental phenomena are difficult to quantify directly. For example, the observed effect may not be the exact placement of a hand, but rather a hand gesture or motion. In such cases, it may not be meaningful to assign a precise uncertainty to the degree or type of gesture.

The key principle is to remain aware of these limitations and to clearly describe them in the dataset documentation. Even if the phenomenon cannot be rigorously quantified, explicitly stating the context, assumptions, and limitations ensures that subsequent users of the data understand its scope and reliability.

\subsection{Expected Results and Required Instruments}

After the quantity to be measured has been established, the expected range - or at least the order of magnitude - of the quantity should be estimated. This is important for selecting appropriate instruments and settings. In many cases, this process is exploratory, involving small probing measurements to gain insight into the expected behavior. For some measurements, this can be straightforward; for example, measuring a DC voltage with a general-purpose voltmeter is usually simple, as these instruments are well-protected, robust, and reliable.

RF measurements, such as those involving spectrum analyzers or vector network analyzers (VNAs), are considerably more complex. These instruments are expensive and sensitive, and misuse can easily damage them. Understanding the expected magnitude of the signal is therefore critical. Spectrum analyzer measurements, for example, are strongly affected by settings such as resolution bandwidth, attenuator setting, pre-amplifier configuration, and reference level. Incorrect settings can alter the internal measurement chain, causing inaccurate readings, increased intermodulation or distortion, spurious harmonics, or a raised noise floor.

By estimating the expected results in advance, the experimenter ensures that instruments are configured appropriately, measurement quality is optimized, and equipment is protected from potential damage. Admittedly, this can be a tedious task at times. Once the approximate magnitude of the quantity is known, the experimenter should consult the instrument manual to verify proper settings or seek guidance from an experienced technician. However, it is important to do preliminary research before asking for help, in order to use the expert’s time efficiently.

Another important reason to approximate the range of the quantity to be measured is that it allows for a preliminary estimation of measurement uncertainty. This early assessment helps determine whether the chosen method and instruments are suitable for the task at hand, or whether an alternative approach is required. By considering the potential uncertainty upfront, the experimenter can avoid investing time in measurements that are unlikely to yield meaningful or reliable results.

As an extension of this, the expected magnitude and required quality of the measurement should guide instrument selection. For example, highly accurate DC voltage measurements may exceed the capabilities of a standard multimeter. Similarly, for AC voltages, if the frequency is in the hundreds of kilohertz or higher, most general-purpose multimeters will not provide reliable readings. While an oscilloscope is very accurate in the time domain, it is not necessarily precise in amplitude, particularly for small signals.

By considering both the expected magnitude and the required measurement fidelity, the experimenter can choose instruments that are suitable for the task and avoid misinterpretation of results caused by inappropriate equipment. Once again, it is strongly recommended to consult an experienced technician before using any non-general-purpose instruments. These devices are often expensive and sensitive, and improper use can result in damage or inaccurate measurements.
\subsection{Selection and Documentation of Measurement Method}

After the quantity to be measured has been established and the appropriate instruments selected, the experimenter should prepare documentation of the measurement principle. In metrology, this is referred to as a **measurement procedure**, or simply a procedure. This is one of the most important tasks of a measurement engineer: documenting the work in sufficient detail so that it can be reproduced with the same results and the same order-of-magnitude uncertainties.

Preparing a measurement procedure is a rigorous task and requires specifying:

\begin{enumerate}
\item What is being measured.
\item Which instruments are used.
\item How stable the system is (e.g., noise, drift, environmental effects).
\item The estimated measurement uncertainty.
\item How everything is connected (a diagram or picture guide is highly recommended).
\end{enumerate}

This is non-negotiable: if data collection is to be performed, the engineer \textbf{must} present a measurement procedure to both the supervisor and laboratory technicians. This ensures that the correct method is used and prevents damage to instruments or personnel.

It should be noted that in experimental work, the requirements for uncertainty estimation may not be as stringent as in industry. However, the experimenter should always remain aware of potential sources of uncertainty and make rough estimates to ensure that the measured quantity is not overwhelmed or obscured by unknown factors.

The procedure should also describe how data is recorded and documented, whether it is collected manually or through an automated system. If automation is used, the code or scripts responsible for data acquisition should be stored alongside the procedure to ensure reproducibility and traceability.


A separate guide will provide more detailed instructions on designing an appropriate measurement procedure. This current section focuses on the principles and documentation requirements, while the other guide offers practical templates and step-by-step guidance.

It should be noted that the measurement procedure can apply to any system, even something as simple as collecting data using an Arduino and its ADC. In such cases, the procedure should describe how the device is used, how it has been verified or calibrated, and any relevant settings or connections. The purpose of the procedure is to ensure that the measurements can always be reproduced and, if necessary, traced back for verification or analysis.

\subsection{Calibration}

Calibration is a cornerstone of reliable measurements. Instruments used for data collection should always be either **calibrated** or **verified against a trusted, calibrated reference**.

In practice, not all laboratory instruments are externally calibrated, as accredited traceable calibration can be expensive. For example, calibrating a common handheld multimeter can cost more than the instrument itself, while calibration of a vector network analyzer (VNA) can easily range from thousands to tens of thousands of euros, often performed annually.

This underscores the principle that expensive instruments must be handled with care: even a minor accident, such as dropping a VNA calibration kit, can necessitate recalibration, resulting in significant cost.

It is therefore the experimenter’s responsibility to ensure that instruments are either properly calibrated or verified against a known reference. This ensures that the measured data reflect the physical quantity of interest, rather than artifacts arising from inaccurate or malfunctioning instruments.
